\section{What a construction must break, and how Markman breaks it}
\label{sec:escape}

\Cref{thm:character} says that a $\Kx$-eigenclass for the norm character has no
component along the Weil line.%
\footnote{The three hypotheses below are not independent in spirit. A construction that imports its input from projective space will normally be equivariant for whatever acts on the input, and an equivariant construction will normally return a single eigenclass, so the three tend to hold or fail together. What \Cref{prop:whichescaped} records is that Markman's construction fails all three at once, and that no two of them can be dropped in isolation.} It is natural to ask what the contrapositive
demands of a construction that does reach the Weil line, and the answer is
sharper than one might expect: the demand is not on the size of the object, nor
on the ambient geometry, nor on any dimension count, but on the way the object
transforms. This section states that demand, and then shows that the one
construction known to meet it, Markman's, meets it in a specific and
reproducible way. The section is the most useful one in the paper for anybody
who wants to build something new, so we have tried to make the requirement
concrete rather than merely negative.

\subsection{The orbit criterion}

\begin{proposition}[Weil classes generate a plane]\label{prop:orbitplane}
Let $(A,\iota,\eta)$ be of $(K,-1,n)$-Weil type and let $\alpha\in\HW(A)$ be
nonzero. Then
\[
  \mathrm{span}_{\QQ}\bigl\{\iota(\tau)^{*}\alpha\;:\;\tau\in\Kx\bigr\}
  \;=\;\HW(A),
\]
a two-dimensional $\QQ$-vector space, and as a representation of $\Kx$ it is
irreducible over $\QQ$.
\end{proposition}

\begin{proof}
By \Cref{prop:weilline}(i) the space $\HW(A)$ is free of rank one over $K$, and
$\iota(\tau)^{*}$ acts on it as multiplication by the scalar $\tau^{2n}\in K$.
Hence the span in question is $S\cdot\alpha$, where
$S=\mathrm{span}_{\QQ}\{\tau^{2n}:\tau\in\Kx\}\subseteq K$. Taking $\tau=1$
shows $1\in S$. By the argument in the proof of \Cref{lem:charsdiffer} the map
$c(\tau)=\tau/\bbar{\tau}$ has infinite image while $\mu(K)$ is finite, so there
is $\tau_{0}$ with $(\tau_{0}/\bbar{\tau_{0}})^{2n}\ne1$, that is with
$\tau_{0}^{2n}\ne\bbar{\tau_{0}}^{\,2n}$, that is with $\tau_{0}^{2n}\notin\QQ$.
Then $S\supseteq\QQ+\QQ\tau_{0}^{2n}=K$, so $S=K$ and the span is
$K\alpha=\HW(A)$.

Irreducibility: a proper nonzero $\Kx$-stable $\QQ$-subspace would be a line,
on which $\Kx$ would act by a $\QQ$-valued character, forcing
$\tau^{2n}\in\QQ$ for all $\tau$, which we have just excluded.
\end{proof}

\begin{corollary}[The orbit criterion]\label{cor:orbitcriterion}
Let $\cE$ be an object of $D^{b}(A)$, or $Z$ a cycle in $\CH^{n}(A)_{\QQ}$,
whose class $\ch_{n}(\cE)$, respectively $\cl(Z)$, is a nonzero Weil class.
Then there is no function $\chi\colon\Kx\to\QQ$ with
\[
  \iota(\tau)^{*}\ch_{n}(\cE)=\chi(\tau)\,\ch_{n}(\cE)
  \qquad\text{for all }\tau\in\Kx .
\]
Equivalently, the $\Kx$-orbit of the class spans a plane and not a line.
\end{corollary}

\begin{proof}
Such a $\chi$ would make $\QQ\cdot\ch_{n}(\cE)$ a $\Kx$-stable line inside
$\HW(A)$, contradicting \Cref{prop:orbitplane}.
\end{proof}

\Cref{cor:orbitcriterion} is the precise sense in which a construction must
break equivariance. Here is what it forbids. Suppose a
construction assigns to the data $(A,\iota,\eta)$ an object $\cE(A,\iota,\eta)$
in a way that is natural, meaning that an isomorphism of the data induces an
isomorphism of the objects. The endomorphism $\iota(\tau)$ is, for $\tau$ in
the order, an isogeny of $A$ commuting with $\iota$ and multiplying $\eta$ by
$\Nm(\tau)$; naturality then forces
$\iota(\tau)^{*}\cE\cong\cE$ up to the normalisation that the polarisation
change imposes, and the induced relation on $\ch_{n}$ is precisely of the
forbidden form $\chi(\tau)\ch_{n}(\cE)$ with $\chi$ a power of the norm. So:

\begin{quote}
\emph{No construction that is natural in the $K$-action, taken alone, can
produce a Weil class. The construction must depend on a choice that
$\Gal(K/\QQ)$ moves.}
\end{quote}

The rest of this section shows what such a choice looks like in practice.

\subsection{The two conjugate isotropic subspaces}

The choice that $\Gal(K/\QQ)$ moves is visible already in \Cref{sec:characters}
and we name it here. The splitting
$V_{\CC}=V_{+}\oplus V_{-}$ of \eqref{eq:Vsplit} is defined over $K$ and not
over $\QQ$: conjugation interchanges the two summands. Consequently
$\bigwedge^{2n}V_{+}$ and $\bigwedge^{2n}V_{-}$ are individually defined over
$K$, and only their sum is rational. A class in $\HW(A)$ is of the form
$w+\bbar{w}$ with $w\in\bigwedge^{2n}V_{+}$, and the two summands are the two
conjugate halves.

\begin{definition}[Conjugate pair data]\label{def:conjpair}
By \emph{conjugate pair data} for $(A,\iota,\eta)$ we mean an ordered choice of
one of the two summands $V_{\pm}$, equivalently an ordered choice of one of the
two embeddings $K\hookrightarrow\CC$, equivalently a choice of one of the two
points of a $\Gal(K/\QQ)$-orbit of size two attached to the $K$-structure.
\end{definition}

\begin{proposition}[Corestriction produces the plane]\label{prop:corestriction}
Let $w\in\bigwedge^{2n}V_{+}$ be nonzero. Then $w+\bbar{w}$ and
$\sqrt{-d}\,(w-\bbar{w})$ lie in $\HW(A)$ and span it over $\QQ$. In
particular, a rational class in the Weil line is exactly a
$\Gal(K/\QQ)$-corestriction of a class in one conjugate half, and any
construction producing such a class must produce, in effect, a pair of
conjugate inputs and combine them.
\end{proposition}

\begin{proof}
Both displayed elements are fixed by conjugation, hence rational, and both lie
in $\bigwedge^{2n}V_{+}\oplus\bigwedge^{2n}V_{-}$ by construction, hence in
$\HW(A)$ by \Cref{def:weilline}. They are linearly independent over $\QQ$
because $w$ and $\bbar{w}$ are linearly independent over $\CC$, and $\HW(A)$ has
dimension two by \Cref{prop:weilline}(i).
\end{proof}

This is the design principle. A construction reaching the Weil line has to be,
in a suitable sense, \emph{bilinear in the two conjugate halves}: it must take
one input attached to $V_{+}$ and one attached to $V_{-}$ and pair them. A
construction fed a single $K$-invariant input cannot do this, which is
\Cref{cor:orbitcriterion} again.

\subsection{How Markman's construction meets the criterion}

Markman's work \cite{Mar25a,Mar25b} is the only construction we know of that
produces Weil classes, and it meets the criterion exactly in the manner just
described. We set out the relevant structure; the details are in
\cite{Mar25b}, and we claim no part of them.

The construction does not begin on $A$ or on $\Av$. It begins on an abelian
$n$-fold $X$, that is, on a variety of half the dimension of $A$. Inside
$H^{\mathrm{ev}}(X,\QQ)$ one selects a plane $P$ spanned by Hodge classes, and
one requires that $P$ meet the spinor variety of even pure spinors in two
points. Those two points are conjugate over the imaginary quadratic field $K$;
they correspond to two maximal isotropic subspaces $W_{1}$ and $W_{2}$ that are
interchanged by $\Gal(K/\QQ)$. This is the conjugate pair data of
\Cref{def:conjpair}, realised concretely.

To each of the two points one attaches a \emph{secant sheaf} on $X$: a coherent
sheaf whose Chern character is the corresponding point of $P$. In the worked
case where $X$ is the Jacobian of a curve of genus three, the secant sheaf is
$\mathcal{I}_{\cup C_{i}}(\Theta)$, the twist by the theta bundle of the ideal
sheaf of a union of translates of the Abel--Jacobi image of the curve. Call the
two sheaves $\cF_{1}$ and $\cF_{2}$.

The abelian variety of Weil type is then $A=X\times\widehat{X}$, of dimension
$2n$, with its $K$-action and polarisation inherited from $P$. The object that
carries the Weil class is
\begin{equation}\label{eq:markmanobject}
  E=\Phi\bigl(\cF_{1}^{\vee}\boxtimes\cF_{2}\bigr),
\end{equation}
where $\Phi\colon D^{b}(X\times X)\to D^{b}(X\times\widehat{X})$ is Orlov's
derived equivalence. One then shows that the characteristic class $\kappa(E)$
remains of Hodge type under every deformation of $X\times\widehat{X}$ as a
polarised abelian $2n$-fold of Weil type, which is what makes the class
algebraic on the whole family. The conclusion in \cite{Mar25b} is the
algebraicity of the Weil classes for abelian sixfolds of Weil type of
discriminant $-1$, and, through Schoen's degeneration argument, for all abelian
fourfolds of Weil type, whence the Hodge conjecture for abelian fourfolds.

\begin{proposition}[Which hypotheses are escaped]\label{prop:whichescaped}
The object \eqref{eq:markmanobject} satisfies none of the hypotheses of
\Cref{thm:divisor} and does not satisfy the hypothesis of
\Cref{thm:equivsuffices}. Specifically:
\begin{enumerate}[label=\textup{(\roman*)}]
\item \emph{The source is not a projective space.} The input
  $\cF_{1}^{\vee}\boxtimes\cF_{2}$ lives on $X\times X$, an abelian variety.
  The proof of \Cref{thm:divisor} used that $H^{*}(\PP^{N-1},\QQ)$ is generated
  in degree two, so that $\ch$ of any object is a polynomial in the hyperplane
  class. The cohomology of $X\times X$ is not generated in degree two: it
  contains K\"unneth classes of every degree, and $\ch(\cF_{1}^{\vee}\boxtimes\cF_{2})$
  is not confined to any divisor subring.
\item \emph{The transform is not $\Phi_{\cP}$ on $A\times\Av$.} The relevant
  equivalence is between $D^{b}(X\times X)$ and $D^{b}(X\times\widehat{X})$,
  and \Cref{thm:fmtheta}, which computes the transform with Poincar\'e kernel
  on powers of the polarisation of $A$, says nothing about it.
\item \emph{The construction is not $K$-equivariant.} The two tensor factors
  of \eqref{eq:markmanobject} carry the two conjugate points of
  $P\cap(\text{spinor variety})$, one each. Interchanging them is the action of
  the nontrivial element of $\Gal(K/\QQ)$, and it exchanges $E$ with a
  different object. Hence $\ch_{n}(E)$ is not a $\Kx$-eigenclass for the norm,
  and \Cref{cor:orbitcriterion} is satisfied rather than violated.
\end{enumerate}
\end{proposition}

\begin{proof}
(i) and (ii) are a comparison of hypotheses and require no argument beyond
reading \eqref{eq:markmanobject}. For (iii), if $\ch_{n}(E)$ were an eigenclass
for the norm then by \Cref{cor:equivariant} it would have zero component along
the Weil line, whereas Markman establishes that it spans one \cite{Mar25b}.
\end{proof}

\begin{remark}[The lesson for the secant idea]\label{rem:secantlesson}
We are explicit here, because the point is easy to miss and is the reason
this section exists. The idea of using secant geometry to produce Weil classes
is not what fails. Markman's construction is a secant construction, and it
works. What fails is siting the secant geometry on the dual $\Av$ of the
$2n$-fold inside an ambient projective space, and transporting by the
Poincar\'e kernel. That siting has three defects at once, by
\Cref{prop:whichescaped}: it puts the input where the cohomology is generated
by a divisor, it uses the transform for which \Cref{thm:fmtheta} applies, and
it feeds the construction $K$-invariant data, since a secant variety of a
$K$-invariant subvariety is $K$-invariant. Moving the secant geometry down to
an abelian $n$-fold repairs all three, and it does so because the halving of
the dimension is what makes room for the two conjugate inputs.
\end{remark}

\subsection{What would be needed above the known range}

We can now state precisely what a construction for the cases that remain open
would have to supply. The known range is $n=2$, all discriminants, and $n=3$
with discriminant $-1$; see \cite{Mar25b} and, for a different route at
discriminant $1$, \cite{FF26}.

\begin{question}\label{q:higher}
Fix $n\ge4$ and an imaginary quadratic field $K$. Does there exist an abelian
$n$-fold $X$, a plane $P\subset H^{\mathrm{ev}}(X,\QQ)$ spanned by Hodge
classes meeting the spinor variety in two points conjugate over $K$, secant
sheaves $\cF_{1},\cF_{2}$ on $X$ realising those two points, and a deformation
argument showing that $\kappa\bigl(\Phi(\cF_{1}^{\vee}\boxtimes\cF_{2})\bigr)$
remains of Hodge type across the moduli of polarised abelian $2n$-folds of
Weil type with multiplication by $K$?
\end{question}

Each of the four ingredients is a concrete demand and none is obviously out of
reach; what is absent at present is the analogue of the genus three Jacobian,
namely a family of abelian $n$-folds in which the plane $P$ and the two secant
sheaves can be written down. By \Cref{cor:orbitcriterion} any successful answer
must, however it is phrased, be bilinear in a pair of conjugate inputs, because
that is the only way a rational class can acquire a $\Kx$-orbit of dimension
two. We offer \Cref{q:higher} as the sharpest form we can give of what is left
on the abelian side, and we note that it is a question about writing down one
sheaf on one abelian variety, not a question about a general mechanism.

\begin{figure}[htbp]
\centering
  \includegraphics[width=0.97\textwidth]{fig_escape}
\caption{The eight positions a secant construction can occupy, drawn as the
Boolean lattice of subsets of the three hypotheses isolated in
\Cref{prop:whichescaped}: (H1) the Chern character of the input lies in the
subring generated by divisor classes, (H2) the output is a $\Kx$-eigenclass,
(H3) the output spans a single line rather than a plane. A construction
satisfying all three sits at the top vertex, where \Cref{thm:divisor} and
\Cref{thm:character} both apply; the shaded face is where (H1) alone already
binds. Markman's construction sits at the bottom vertex, satisfying none of
the three, which is why nothing here bears on it, and the Weil line needs a
plane by \Cref{prop:orbitplane}.}
\label{fig:escape}
\end{figure}
